% !TEX program = lualatex
\documentclass[11pt,a4paper]{article}
\usepackage{luatexja}
\usepackage{amsmath,amssymb,amsthm,amsfonts}
\usepackage{geometry}
\geometry{margin=1in}

% 定理環境の設定
\theoremstyle{definition}
\newtheorem{definition}{定義}[section]
\newtheorem{theorem}{定理}[section]
\newtheorem{remark}{備考}[section]

\title{真理保存性の不完全性定理に基づく数論的困難性の解明：\\コラッツ予想における「1のフラクタル」と次元上昇の必然性}
\author{千葉将臣}
\date{2026-04-29}

\begin{document}

\maketitle

\begin{abstract}
本稿は、先行研究で定式化された「真理保存性の不完全性定理」を自然数論における未解決問題、特にコラッツ予想（$3n+1$問題）に適用する。古典的算術体系が天下りに要請する「1の自己同一性（素因数分解の一意性）」が、乗法と加法の干渉によって外部へ無限に伸展する「1のフラクタル」を生成することを示す。本定理の帰結として、このフラクタルな伸展の大域的収束性を、自己同一性を保存したままの2次元的推論によって証明することは構造上不可能であることを論証し、位相幾何学的な縮小写像（0のフラクタル）へと問題を再翻訳する「界論的次元上昇」の必然性を提示する。
\end{abstract}

\section{導入}
数学基礎論において、体系が自己の無矛盾性や不変量を要請する限り、その内部からは証明不能な命題が必然的に生じる（真理保存性の不完全性定理）。本稿では、このメタ定理が単なる自己言及論理の限界に留まらず、コラッツ予想に代表される現実の数論的未解決問題が抱える本質的困難性の原因であることを明らかにする。

\section{コラッツ力学系における「1のフラクタル」の生成}

\begin{definition}[自己同一性の伸展としての「1のフラクタル」]
自然数論の基盤は、単位元「1」の自己同一性と、それに基づく素因数分解の一意性にある。コラッツ操作における偶数ルール（乗法的縮小）と奇数ルール（加法と乗法の混合：$3n+1$）は、互いに素である「2」と「3」の干渉を引き起こす。この干渉は、要素が自己同一性を維持しようとする反発力として働き、系全体を「1」を起点として位相空間の外部へ無限に分岐・伸展させる構造（コラッツ・ツリー）を生成する。本稿ではこれを「1のフラクタル」と定義する。
\end{definition}

古典的算術（2次元的観測者）は、対象の自己同一性（$A=A$）を真理保存の前提とするため、個々の自然数 $n$ の軌道を決定論的なステップとして追跡する手法をとる。

\section{真理保存性の不完全性定理の適用}

先行研究における「真理保存性の不完全性定理」を本系に適用すると、以下の構造的限界が導出される。

\begin{theorem}[コラッツ予想の内部証明不可能性]
自然数論の体系 $T$ は「自己同一性の保存」を前提として稼働する。体系 $T$ は、その前提自体が生み出した無限の外部伸展（1のフラクタル）の全体像や、その限界（大域的収束性）を、$T$ 内部の局所的な演繹操作によって証明することはできない。
\end{theorem}

証明の構造的理由は明快である。「証明」とは本質的に多様な状態を一つの結論へと圧縮・収束させる行為である。しかし、1のフラクタルは定義上、系が真理（自己同一性）を保存しようとする際の「自己参照の閉じ損ね」によって外部へと発散し続ける軌道である。系が自身の足場（真理保存性）を疑わずに、自らが生み出した無限の発散のすべてをしらみつぶしに検証し、有限のステップで閉包を宣言することは、論理的に循環あるいは発散する。これが、コラッツ予想が既存の数学的プロトコルにおいて未知であり続ける最大の構造的要因である。

\section{次元上昇による「0のフラクタル」への反転}

古典的体系におけるこの証明不可能性（影）を突破するためには、真理保存性の前提（個別の数の自己同一性）を一旦放棄し、メタ階層へと観測者の次元を上昇させる必要がある。

\begin{remark}[界論的アプローチの必然性]
観測空間を高階位相空間（8次元観測空間）へと拡張し、無限の自然数の束全体を一つの測度空間として対象化する。これにより、外部へ伸展する「1のフラクタル」は、系全体のエントロピー縮減（縮小写像）に伴う「0のフラクタル（大域的アトラクタへの変位レトラクト）」へと幾何学的に反転される。
\end{remark}

個別の同一性に執着する限り証明不能であった問題は、空間全体が位相的特異点へと収束・吸収されるプロセスとして再定式化される。このとき初めて、証明は「特定の基本群（例：$H^1 \cong \mathbb{Z}$）の導出」という代数トポロジーの計算対象として既存数学と再接続される。

\section{結論}
コラッツ予想の難解さは、計算の複雑さにあるのではなく、古典数学が天下りに要請する「自己同一性の真理保存」がもたらす構造的な不完全性にある。本問題の解決には、真理保存性の呪縛を脱し、対象を大域的な位相的収束へと変換する次元数学的アプローチが不可避である。

\end{document}
